\chapter{第三单元}

\Needspace{12\baselineskip}

\section{第三单元素养评价（A）}

\Needspace{8\baselineskip}

\pdfbookmark[2]{原书第17页}{orig-17}

\subsection*{原书第17页}

% p017_s01
\BookQuestion{一、1（全部空）}
妈妈花45元买了3千克草莓，每千克草莓多少元。“45元”叫作（\hspace{1.8em}），“3千克”叫作（\hspace{1.8em}），每千克草莓的价钱叫作（\hspace{1.8em}）。数量关系式：\BookBlank{}，列式计算：\BookBlank{}。\par
\begin{BookAnswer}
\(45\)元是买草莓的\textbf{总价}，\allowbreak{}\(3\)千克是\textbf{数量}，每千克的价钱叫\textbf{单价}。数量关系为“单价\(=\)总价\(\div\)数量”。\[\begin{aligned}&45\div3\\&=15\text{(元/千克)}\end{aligned}\]所以草莓每千克\(15\)元。\par 用“总价\(=\)单价\(\times\)数量”检验：\(15\times3=45\text{(元)}\)，与妈妈花的钱一致。
\end{BookAnswer}

% p017_s02
\BookQuestion{一、2（牛奶）}
将表格补充完整：牛奶65元/箱，总价195元，数量应填（\hspace{1.8em}）。\par
\par\smallskip\Needspace{4\baselineskip}\begingroup\small\setlength{\tabcolsep}{3pt}
\noindent\begin{tabularx}{\linewidth}{|*{4}{>{\centering\arraybackslash}X|}}\hline
物品 & 单价 & 数量 & 总价\\\hline
牛奶 & 65元/箱 & （    ） & 195元\\\hline
鸡蛋 & （    ） & 6盒 & 108元\\\hline
奶茶 & 12元/杯 & 5杯 & （    ）\\\hline
\end{tabularx}\par
\endgroup\smallskip
\begin{BookAnswer}
已知总价和单价，求数量：\(195\div65=3\text{(箱)}\)。数量栏填\(3\)箱。
\end{BookAnswer}

% p017_s03
\BookQuestion{一、2（鸡蛋）}
将表格补充完整：鸡蛋6盒，总价108元，单价应填（\hspace{1.8em}）。\par
\par\smallskip\Needspace{4\baselineskip}\begingroup\small\setlength{\tabcolsep}{3pt}
\noindent\begin{tabularx}{\linewidth}{|*{4}{>{\centering\arraybackslash}X|}}\hline
物品 & 单价 & 数量 & 总价\\\hline
牛奶 & 65元/箱 & （    ） & 195元\\\hline
鸡蛋 & （    ） & 6盒 & 108元\\\hline
奶茶 & 12元/杯 & 5杯 & （    ）\\\hline
\end{tabularx}\par
\endgroup\smallskip
\begin{BookAnswer}
已知总价和数量，求单价：\(108\div6=18\text{(元/盒)}\)。单价栏填\(18\)元/盒。
\end{BookAnswer}

% p017_s04
\BookQuestion{一、2（奶茶）}
将表格补充完整：奶茶12元/杯，数量5杯，总价应填（\hspace{1.8em}）。\par
\par\smallskip\Needspace{4\baselineskip}\begingroup\small\setlength{\tabcolsep}{3pt}
\noindent\begin{tabularx}{\linewidth}{|*{4}{>{\centering\arraybackslash}X|}}\hline
物品 & 单价 & 数量 & 总价\\\hline
牛奶 & 65元/箱 & （    ） & 195元\\\hline
鸡蛋 & （    ） & 6盒 & 108元\\\hline
奶茶 & 12元/杯 & 5杯 & （    ）\\\hline
\end{tabularx}\par
\endgroup\smallskip
\begin{BookAnswer}
已知单价和数量，求总价：\(12\times5=60\text{(元)}\)。总价栏填\(60\)元。乘法按“每杯价格\(\times\)杯数”列式。
\end{BookAnswer}

% p017_s05
\BookQuestion{一、3（全部空）}
高铁列车每小时约行驶305千米，这是这列高铁列车的（\hspace{1.8em}），可以写成（\hspace{1.8em}），读作（\hspace{1.8em}），这列高铁列车从甲地到乙地需要8小时，甲地到乙地全程约有（\hspace{1.8em}）千米。\par
\begin{BookAnswer}
每小时行\(305\)千米表示\textbf{速度}，写作\(305\text{千米/时}\)，读作“\(305\)千米每时”。\par 由“路程\(=\)速度\(\times\)时间”：\[\begin{aligned}&305\times8\\&=2440\text{(千米)}\end{aligned}\]路程空填\(2440\)。
\end{BookAnswer}

% p017_s06
\BookQuestion{一、4（两个空）}
王老师每分钟大约走100米，（\hspace{1.8em}）分钟大约走1千米。数量关系式：\BookBlank{}。\par
\begin{BookAnswer}
先统一长度单位：\[1\text{千米}=1000\text{米}\]所需时间\(1000\div100=10\text{(分)}\)，数值空填\(10\)；关系式填“时间\(=\)路程\(\div\)速度”。
\end{BookAnswer}

% p017_s07
\BookQuestion{一、5}
每根跳绳8元，老师用245元最多能买（\hspace{1.8em}）根跳绳。\par
\begin{BookAnswer}
\(245=8\times30+5\)，买\(30\)根后剩\(5\)元，不够再买一根；又有\(8\times31=248>245\)。所以最多买\(30\)根。
\end{BookAnswer}

% p017_s08
\BookQuestion{一、6}
大货车5小时行驶了360千米，冷藏车4小时行驶了320千米。两辆车相比，（\hspace{1.8em}）的速度更快。\par
\begin{BookAnswer}
比较速度而不是只比较路程。大货车\(360\div5=72\text{(千米/时)}\)，冷藏车\(320\div4=80\text{(千米/时)}\)。\(80>72\)，填“冷藏车”。
\end{BookAnswer}

% p017_s09
\BookQuestion{一、7}
万达影城的2号厅播放一部电影，该厅共有102个座位，全部坐满，每人票价38元，则该厅共能收入（\hspace{1.8em}）元。\par
\begin{BookAnswer}
座位全部坐满，收入\(=\)每张票价\(\times\)售票张数。\[\begin{aligned}&38\times102\\&=38\times(100+2)\\&=3800+76\\&=3876\text{(元)}\end{aligned}\]填\(3876\)。
\end{BookAnswer}

% p017_s10
\BookQuestion{一、8（跨页）}
新疆棉日照充足、纤维细长、柔软透气、品质优良，李阿姨用288元买了一批新疆棉毛巾，一共买了8盒，每盒有3条，每盒毛巾单价是（\hspace{1.8em}）元，平均每条毛巾（\hspace{1.8em}）元。（后半题续于原书第18页。）\par
\begin{BookAnswer}
题目续到下一页：\(288\)元买\(8\)盒，每盒\(3\)条。每盒单价\(288\div8=36\)元，每条单价\(36\div3=12\)元。两空分别为\(36,12\)，完整脱式与另一方法见下一页。
\end{BookAnswer}

\Needspace{8\baselineskip}

\pdfbookmark[2]{原书第18页}{orig-18}

\subsection*{原书第18页}

% p018_s01
\BookQuestion{一、8（续）}
（接原书第17页第8题）李阿姨用288元买了8盒新疆棉毛巾，每盒有3条，每盒毛巾单价是（\hspace{1.8em}）元，平均每条毛巾（\hspace{1.8em}）元。\par
\begin{BookAnswer}
两个空依次填\(36,12\)。原书参考把每条写作\(10\)元，与题面\(288\)元、\(8\)盒、每盒\(3\)条不符，应按题面核算。\par \noindent\begin{minipage}{\linewidth}\textbf{方法一：先求每盒，再求每条}\par \[\begin{aligned}&288\div8\\&=36\text{(元)}\end{aligned}\]\[\begin{aligned}&288\div8\div3\\&=36\div3\\&=12\text{(元)}\end{aligned}\]答：每盒\(36\)元，每条\(12\)元。\end{minipage}\par \noindent\begin{minipage}{\linewidth}\textbf{方法二：先求总条数，再求每条}\par \[\begin{aligned}&288\div(3\times8)\\&=288\div24\\&=12\text{(元)}\end{aligned}\]每盒有\(3\)条，所以每盒价格为：\[\begin{aligned}&12\times3\\&=36\text{(元)}\end{aligned}\]答：每盒\(36\)元，每条\(12\)元。\end{minipage}\par \par 这里的两个单价分别按“元/盒”和“元/条”计，不能混用。检验\(12\times3\times8=288\text{(元)}\)，与总价相符；也有\(36\div3=12\)，每条价应小于每盒价。
\end{BookAnswer}

% p018_s02
\BookQuestion{一、9（1）}
国庆出游，王叔叔开车带家人自驾游，其中高速路段车速平均每小时106千米，全程累计开了22小时，一共开了多少千米？（1）解决这个问题所依据的数量关系式是（\hspace{1.8em}）。\par
\begin{BookAnswer}
数量关系式为“路程\(=\)速度\(\times\)时间”。\par 原文只说明“高速路段”的平均速度是\(106\)千米/时，却给出“全程”的\(22\)小时。只有当全程这\(22\)小时都按此平均速度计算时，才有\[\begin{aligned}&106\times22\\={}&2332\text{(千米)}\end{aligned}\]若全程还包含其他路段且其速度未知，则不能仅凭所给条件唯一求全程路程。
\end{BookAnswer}

% p018_s03
\BookQuestion{一、9（2）（两个空）}
国庆出游，王叔叔开车带家人自驾游，其中高速路段车速平均每小时106千米，全程累计开了22小时，一共开了多少千米？（2）右图竖式中，第②步计算是\(106\times\)（\hspace{1.8em}），表示的是（\hspace{1.8em}）小时行驶的路程。\par
\BookFigure{figures/p018_s03-01.tex}
\begin{BookAnswer}
把\(22\)拆成\(20+2\)，则\(106\times22=106\times20+106\times2\)。竖式的两行部分积正是这两个乘积。\par 竖式第二个部分积中的\(2\)在十位，实际表示\(20\)。第②步算的是\(106\times20=2120\)，表示行驶\(20\)小时的路程；两空都填\(20\)，不能填\(2\)。
\end{BookAnswer}

% p018_s04
\BookQuestion{二、1}
小华步行去商店，路上用了10分钟；买了3本练习本，花了15元。这里的10分钟、15元、3本分别表示的是（\hspace{1.8em}）。\par
\noindent A.路程，总价，数量\par
\noindent B.时间，总价，单价\par
\noindent C.时间，总价，数量\par
\noindent D.速度，单价，数量\par
\begin{BookAnswer}
选\(\mathrm C\)。\(10\)分钟表示\textbf{时间}，\allowbreak{}\(15\)元表示三本练习本的\textbf{总价}，\allowbreak{}\(3\)本表示\textbf{数量}。
\end{BookAnswer}

% p018_s05
\BookQuestion{二、2（两个空）}
每瓶矿泉水3元，买12瓶一共花了多少钱？用（\hspace{1.8em}）计算；一名外卖骑手，5分钟骑行1400米，求速度，用（\hspace{1.8em}）计算。\par
\noindent A.加法\par
\noindent B.乘法\par
\noindent C.除法\par
\noindent D.减法\par
\begin{BookAnswer}
第一空选\(\mathrm B\)：单价乘数量，\allowbreak{}\(3\times12=36\text{(元)}\)，用乘法。\par 第二空选\(\mathrm C\)：速度等于路程除以时间，\allowbreak{}\(1400\div5=280\text{(米/分)}\)，用除法。
\end{BookAnswer}

% p018_s06
\BookQuestion{二、3}
如图，一列货车从甲地开往乙地用了4个小时，货车的速度是（\hspace{1.8em}）。\par
\noindent A.70时\par
\noindent B.70千米\par
\noindent C.70千米/时\par
\noindent D.280千米/时\par
\BookFigure{figures/p018_s06-01.tex}
\begin{BookAnswer}
选\(\mathrm C\)。\(280\div4=70\)，速度必须同时表示“路程和时间”的比，单位是“千米/时”，不是单独的时或千米。
\end{BookAnswer}

% p018_s07
\BookQuestion{二、4}
下面问题中，适合用“速度\(\times\)时间”来解决的是（\hspace{1.8em}）。\par
\noindent A.求每支钢笔的钱数\par
\noindent B.求汽车5小时行驶的千米数\par
\noindent C.求小明每分钟步行的米数\par
\noindent D.计算长方形的周长\par
\begin{BookAnswer}
选\(\mathrm B\)。已知汽车速度和行驶\(5\)小时的时间，求行驶路程，要用“速度\(\times\)时间”。\(\mathrm C\)是已知路程和时间求速度，应用除法。
\end{BookAnswer}

% p018_s08
\BookQuestion{二、5}
这个公路指示牌告诉我们当前位置距青岛的（\hspace{1.8em}）还有203千米。\par
\noindent A.时间\par
\noindent B.速度\par
\noindent C.路程\par
\noindent D.数量\par
\par\smallskip\Needspace{4\baselineskip}\begingroup\small\setlength{\tabcolsep}{3pt}
\noindent\begin{tabularx}{\linewidth}{|*{2}{>{\centering\arraybackslash}X|}}\hline
目的地 & 距离\\\hline
青岛 & 203千米\\\hline
\end{tabularx}\par
\endgroup\smallskip
\begin{BookAnswer}
选\(\mathrm C\)。\(203\)千米说明当前位置到青岛还有多远，表示\textbf{路程}；“千米”不是时间或速度的单位。
\end{BookAnswer}

% p018_s09
\BookQuestion{二、6（A、C都正确）}
张平用18元买了6本笔记本，小华买这样的4本笔记本花了多少钱？列式正确的是（\hspace{1.8em}）。\par
\noindent A.\(18\times4\div6\)\par
\noindent B.\(18\times6\div4\)\par
\noindent C.\(18\div6\times4\)\par
\noindent D.\(18\div6+4\)\par
\begin{BookAnswer}
本题按单选形式设置不恰当，\allowbreak{}\(\mathrm A,\mathrm C\)都成立。\par \par\noindent\begin{minipage}{\linewidth}\textbf{C：先求单价，再求四本的钱}\par \[\begin{aligned}&18\div6\times4\\&=3\times4\\&=12\text{(元)}\end{aligned}\]\end{minipage}\par \par \par\noindent\begin{minipage}{\linewidth}\textbf{A：先求四组的钱，再平均分}\par \[\begin{aligned}&18\times4\div6\\&=72\div6\\&=12\text{(元)}\end{aligned}\]\end{minipage}\par \(4\)组、每组\(6\)本共\(24\)本，把总价平均分成\(6\)份，每份正好是\(4\)本的钱。两式都得到\(12\)元；不能因为\(\mathrm A\)先乘后除就判错。\par 两种列式还可用相同单价核对：\(12\div4=18\div6=3\text{(元/本)}\)，所以\(A,C\)都保持了原题单价。
\end{BookAnswer}

\Needspace{8\baselineskip}

\pdfbookmark[2]{原书第19页}{orig-19}

\subsection*{原书第19页}

% p019_s01
\BookQuestion{三、1（1）}
直接写得数：\(125\times8=\)。\par
\begin{BookAnswer}
把125看作1000的八分之一。\[\begin{aligned} &125\times8\\={}&1000\end{aligned}\]
\end{BookAnswer}

% p019_s02
\BookQuestion{三、1（2）}
直接写得数：\(40\times25=\)。\par
\begin{BookAnswer}
先用4与25凑100。\[\begin{aligned} &40\times25\\={}&4\times25\times10\\={}&100\times10\\={}&1000\end{aligned}\]
\end{BookAnswer}

% p019_s03
\BookQuestion{三、1（3）}
直接写得数：\(103-74=\)。\par
\begin{BookAnswer}
按减法计算。\[\begin{aligned} &103-74\\={}&29\end{aligned}\]
\end{BookAnswer}

% p019_s04
\BookQuestion{三、1（4）}
直接写得数：\(0\times75+9=\)。\par
\begin{BookAnswer}
先乘后加；0乘任何数都得0。\[\begin{aligned} &0\times75+9\\={}&0+9\\={}&9\end{aligned}\]
\end{BookAnswer}

% p019_s05
\BookQuestion{三、1（5）}
直接写得数：\(13+57\times2=\)。\par
\begin{BookAnswer}
先乘后加。\[\begin{aligned} &13+57\times2\\={}&13+114\\={}&127\end{aligned}\]
\end{BookAnswer}

% p019_s06
\BookQuestion{三、1（6）}
直接写得数：\(188+423=\)。\par
\begin{BookAnswer}
按数位对齐相加。\[\begin{aligned} &188+423\\={}&611\end{aligned}\]
\end{BookAnswer}

% p019_s07
\BookQuestion{三、1（7）}
直接写得数：\(50\times18=\)。\par
\begin{BookAnswer}
5乘18得90，再乘10。\[\begin{aligned} &50\times18\\={}&900\end{aligned}\]
\end{BookAnswer}

% p019_s08
\BookQuestion{三、1（8）}
直接写得数：\(67\times5=\)。\par
\begin{BookAnswer}
可算67乘10再取一半。\[\begin{aligned} &67\times5\\={}&335\end{aligned}\]
\end{BookAnswer}

% p019_s09
\BookQuestion{三、1（9）}
直接写得数：\(32\div2\div4=\)。\par
\begin{BookAnswer}
同级运算从左往右。\[\begin{aligned} &32\div2\div4\\={}&16\div4\\={}&4\end{aligned}\]
\end{BookAnswer}

% p019_s10
\BookQuestion{三、1（10）}
直接写得数：\(56-34\div2=\)。\par
\begin{BookAnswer}
先除后减。\[\begin{aligned} &56-34\div2\\={}&56-17\\={}&39\end{aligned}\]
\end{BookAnswer}

% p019_s11
\BookQuestion{三、2（1）}
计算下列各题：（1）\(360\div(12+360\div15)\)。\par
\begin{BookAnswer}
先算小括号内的除法。\[\begin{aligned} &360\div(12+360\div15)\\={}&360\div(12+24)\\={}&360\div36\\={}&10\end{aligned}\]
\end{BookAnswer}

% p019_s12
\BookQuestion{三、2（2）}
计算下列各题：（2）\(80\times[(103+97)-131]\)。\par
\begin{BookAnswer}
先小括号，再中括号。\[\begin{aligned} &80\times[(103+97)-131]\\={}&80\times[200-131]\\={}&80\times69\\={}&5520\end{aligned}\]\par\textbf{简便提示（学完运算律后）}\par \[\begin{aligned}&80\times69\\={}&80\times(70-1)\\={}&5600-80\\={}&5520\end{aligned}\]
\end{BookAnswer}

% p019_s13
\BookQuestion{三、2（3）}
计算下列各题：（3）\(680-85\times4\div20\)。\par
\begin{BookAnswer}
乘除同级，从左往右。\[\begin{aligned} &680-85\times4\div20\\={}&680-340\div20\\={}&680-17\\={}&663\end{aligned}\]\par\textbf{简便提示（学完运算律后）}\par \[\begin{aligned}&680-85\times4\div20\\={}&680-85\times4\div(4\times5)\\={}&680-85\div5\\={}&680-17\\={}&663\end{aligned}\]
\end{BookAnswer}

% p019_s14
\BookQuestion{三、2（4）}
计算下列各题：（4）\(396\div(25+11)\times12\)。\par
\begin{BookAnswer}
先算小括号，再依次除、乘。\[\begin{aligned} &396\div(25+11)\times12\\={}&396\div36\times12\\={}&11\times12\\={}&132\end{aligned}\]
\end{BookAnswer}

% p019_s15
\BookQuestion{四、1}
我国自主研发的物流货运无人机，在山区邮政运输中，单次最大载重120千克。照这样计算，35架同款无人机各运送一次，一共可以运送货物多少千克？\par
\begin{BookAnswer}
主解：每架载质量×架数\par \[\begin{aligned} &120\times35\\={}&4200\text{(千克)}\end{aligned}\]答：一共可以运送4200千克货物。\par 同一列式的简算：把35架分成30架和5架\par \[\begin{aligned} &120\times(30+5)\\={}&120\times30+120\times5\\={}&3600+600\\={}&4200\text{(千克)}\end{aligned}\]答：一共可以运送4200千克货物。\par
\end{BookAnswer}

% p019_s16
\BookQuestion{四、2}
学校为校园读书节采购红色经典读本，一共买了25本，总共花费600元，每本红色经典读本的单价是多少元？\par
\begin{BookAnswer}
数量关系：总价＝单价×数量。已知600元买25本，求每本价格，用总价÷数量。\par 主解：总价÷数量\par \[\begin{aligned} &600\div25\\={}&24\text{(元)}\end{aligned}\]答：每本红色经典读本24元。\par 计算补充：用连续除法简算\par \[\begin{aligned} &600\div25\\={}&600\div(5\times5)\\={}&600\div5\div5\\={}&120\div5\\={}&24\text{(元)}\end{aligned}\]答：每本红色经典读本24元。\par 检验：\(24\times25=600\text{(元)}\)，回到已知总价。
\end{BookAnswer}

\Needspace{8\baselineskip}

\pdfbookmark[2]{原书第20页}{orig-20}

\subsection*{原书第20页}

% p020_s01
\BookQuestion{四、3}
社区为孤寡老人采购保暖物资，买18套毛毯，每套45元，付给收银员900元，应找回多少元？\par
\begin{BookAnswer}
方法一：付出的钱−购买18套的钱\par \[\begin{aligned} &900-45\times18\\={}&900-810\\={}&90\text{(元)}\end{aligned}\]答：应找回90元。\par 900元正好能买20套，只买18套，剩2套的钱。\par 方法二：每套价格×少买的套数\par \[\begin{aligned} &45\times(900\div45-18)\\={}&45\times(20-18)\\={}&45\times2\\={}&90\text{(元)}\end{aligned}\]答：应找回90元。\par
\end{BookAnswer}

% p020_s02
\BookQuestion{四、4}
一辆汽车从甲地开往乙地，速度是60千米/时，4小时到达，原路返回时只用了3小时，返回时平均每小时行多少千米？\par
\begin{BookAnswer}
往返的路程相同，先求单程路程，再除以返回时间。\par \textbf{路程÷返回时间}\par \[\begin{aligned} &60\times4\div3\\={}&240\div3\\={}&80\text{(千米/时)}\end{aligned}\]答：返回时平均每小时行80千米。\par \par\textbf{简便提示（学完运算律后）}\par \[\begin{aligned}&60\times4\div3\\={}&60\div3\times4\\={}&20\times4\\={}&80\text{(千米/时)}\end{aligned}\]
\end{BookAnswer}

% p020_s03
\BookQuestion{四、5}
一辆面包车和一辆轿车分别从济南市、威海市同时相对开出，3小时相遇。已知面包车平均每小时行驶80千米，轿车平均每小时行驶100千米，济南市和威海市相距多少千米？（先画图整理条件和问题，再解答。）\par
\begin{BookAnswer}
先核对同时出发、相向行驶、3小时后相遇：两车用时相同，且相遇时没有剩余间距。图中两段分别是两车3小时的路程，两段之和就是两市距离。\par 方法一：速度和×相遇时间\par \[\begin{aligned} &(80+100)\times3\\={}&180\times3\\={}&540\text{(千米)}\end{aligned}\]答：济南市与威海市相距540千米。\par 方法二：分别求路程再相加\par \[\begin{aligned} &80\times3+100\times3\\={}&240+300\\={}&540\text{(千米)}\end{aligned}\]答：济南市与威海市相距540千米。\par 检验：\(540\div3-100=80\text{(千米/时)}\)，可以还原面包车的已知速度。
\BookFigure{figures/p020_s03-01.tex}
\end{BookAnswer}

% p020_s04
\BookQuestion{四、6}
亮亮一家3口计划到海南五日游，现在有两种方案：①组团游：成人2200元/人，18岁以下1800元/人（含餐费和住宿费）。②自由游：飞机票单程900元/人；住宿费每人每天120元；吃饭、交通费每人每天120元。算一算，哪种方案更划算？\par
\begin{BookAnswer}
一家3口按2名成人和1名儿童计算；自由游机票应付往返费用。\par 方法一：先分别合计交通和食宿费用\par 组团游：\[\begin{aligned} &2200\times2+1800\\={}&4400+1800\\={}&6200\text{(元)}\end{aligned}\]自由游：\[\begin{aligned} &900\times2\times3+(120+120)\times5\times3\\={}&1800\times3+240\times5\times3\\={}&5400+1200\times3\\={}&5400+3600\\={}&9000\text{(元)}\end{aligned}\]\(6200<9000\)，且\(9000-6200=2800\text{(元)}\)。答：组团游更划算，少花2800元。\par 方法二：先算自由游每人的费用\par \[\begin{aligned} &900\times2+(120+120)\times5\\={}&1800+240\times5\\={}&1800+1200\\={}&3000\text{(元)}\end{aligned}\]\[\begin{aligned} &3000\times3-6200\\={}&9000-6200\\={}&2800\text{(元)}\end{aligned}\]答：组团游更划算，少花2800元。
\end{BookAnswer}

\Needspace{12\baselineskip}

\section{第三单元素养评价（B）}

\Needspace{8\baselineskip}

\pdfbookmark[2]{原书第21页}{orig-21}

\subsection*{原书第21页}

% p021_s01
\BookQuestion{一、1}
买3支钢笔需要96元，买10支这样的钢笔一共需要（\hspace{1.8em}）元。\par
\begin{BookAnswer}
先求每支钢笔价格，再求10支总价：\[\begin{aligned} &96\div3\times10\\={}&32\times10\\={}&320\text{(元)}\end{aligned}\]
\end{BookAnswer}

% p021_s02
\BookQuestion{一、2（已知条件）}
8千克苹果40元，1千克苹果多少元？这道题是已知（\hspace{1.8em}）和（\hspace{1.8em}），求（\hspace{1.8em}），数量关系式是：\BookBlank{}。\par
\begin{BookAnswer}
40元是总价，8千克是数量，问题求单价。数量关系为：单价＝总价÷数量。
\end{BookAnswer}

% p021_s03
\BookQuestion{一、2（单价和读法）}
8千克苹果40元，苹果的单价是（\hspace{1.8em}），读作（\hspace{1.8em}）。\par
\begin{BookAnswer}
\[\begin{aligned} &40\div8\\={}&5\text{(元/千克)}\end{aligned}\]读作“5元每千克”，表示每千克苹果5元。
\end{BookAnswer}

% p021_s04
\BookQuestion{一、3（汽车）}
一辆汽车每小时行驶72千米，可以写作（\hspace{1.8em}），读作（\hspace{1.8em}）。\par
\begin{BookAnswer}
写作\(72\text{千米/时}\)，读作“72千米每时”；分母“时”表示每1小时。
\end{BookAnswer}

% p021_s05
\BookQuestion{一、3（步行）}
李亮每分钟步行65米，可以写作（\hspace{1.8em}），读作（\hspace{1.8em}）。\par
\begin{BookAnswer}
写作\(65\text{米/分}\)，读作“65米每分”；表示每分钟走65米。
\end{BookAnswer}

% p021_s06
\BookQuestion{一、4（步行速度）}
把表格补充完整：步行路程840米，时间12分，速度（\hspace{1.8em}）。\par
\par\smallskip\Needspace{4\baselineskip}\begingroup\small\setlength{\tabcolsep}{3pt}
\noindent\begin{tabularx}{\linewidth}{|*{5}{>{\centering\arraybackslash}X|}}\hline
项目 & 步行 & 自行车 & 船 & 汽车\\\hline
路程 & 840米 & （    ） & 105千米 & 428千米\\\hline
时间 & 12分 & 5分 & （    ） & 4小时\\\hline
速度 & （    ） & 280米/分 & 35千米/时 & （    ）\\\hline
\end{tabularx}\par
\endgroup\smallskip
\begin{BookAnswer}
速度＝路程÷时间：\[\begin{aligned} &840\div12\\={}&70\text{(米/分)}\end{aligned}\]
\end{BookAnswer}

% p021_s07
\BookQuestion{一、4（骑车路程）}
把表格补充完整：骑自行车速度280米/分，时间5分，路程（\hspace{1.8em}）。\par
\par\smallskip\Needspace{4\baselineskip}\begingroup\small\setlength{\tabcolsep}{3pt}
\noindent\begin{tabularx}{\linewidth}{|*{5}{>{\centering\arraybackslash}X|}}\hline
项目 & 步行 & 自行车 & 船 & 汽车\\\hline
路程 & 840米 & （    ） & 105千米 & 428千米\\\hline
时间 & 12分 & 5分 & （    ） & 4小时\\\hline
速度 & （    ） & 280米/分 & 35千米/时 & （    ）\\\hline
\end{tabularx}\par
\endgroup\smallskip
\begin{BookAnswer}
路程＝速度×时间：\[\begin{aligned} &280\times5\\={}&1400\text{(米)}\end{aligned}\]
\end{BookAnswer}

% p021_s08
\BookQuestion{一、4（船行时间）}
把表格补充完整：船行路程105千米，速度35千米/时，时间（\hspace{1.8em}）。\par
\par\smallskip\Needspace{4\baselineskip}\begingroup\small\setlength{\tabcolsep}{3pt}
\noindent\begin{tabularx}{\linewidth}{|*{5}{>{\centering\arraybackslash}X|}}\hline
项目 & 步行 & 自行车 & 船 & 汽车\\\hline
路程 & 840米 & （    ） & 105千米 & 428千米\\\hline
时间 & 12分 & 5分 & （    ） & 4小时\\\hline
速度 & （    ） & 280米/分 & 35千米/时 & （    ）\\\hline
\end{tabularx}\par
\endgroup\smallskip
\begin{BookAnswer}
时间＝路程÷速度：\[\begin{aligned} &105\div35\\={}&3\text{(小时)}\end{aligned}\]
\end{BookAnswer}

% p021_s09
\BookQuestion{一、4（汽车速度）}
把表格补充完整：汽车路程428千米，时间4小时，速度（\hspace{1.8em}）。\par
\par\smallskip\Needspace{4\baselineskip}\begingroup\small\setlength{\tabcolsep}{3pt}
\noindent\begin{tabularx}{\linewidth}{|*{5}{>{\centering\arraybackslash}X|}}\hline
项目 & 步行 & 自行车 & 船 & 汽车\\\hline
路程 & 840米 & （    ） & 105千米 & 428千米\\\hline
时间 & 12分 & 5分 & （    ） & 4小时\\\hline
速度 & （    ） & 280米/分 & 35千米/时 & （    ）\\\hline
\end{tabularx}\par
\endgroup\smallskip
\begin{BookAnswer}
速度＝路程÷时间：\[\begin{aligned} &428\div4\\={}&107\text{(千米/时)}\end{aligned}\]
\end{BookAnswer}

% p021_s10
\BookQuestion{一、5}
由中国自主研发、世界领先的新型复兴号高速综合检测列车上线运行后，在世界上首次成功实现明线单列时速435千米。照这样的速度，新型复兴号高速综合检测列车12小时能行驶（\hspace{1.8em}）千米，用到的数量关系是：\BookBlank{}。\par
\begin{BookAnswer}
数量关系：路程＝速度×时间。\[\begin{aligned} &435\times12\\={}&435\times(10+2)\\={}&4350+870\\={}&5220\text{(千米)}\end{aligned}\]
\end{BookAnswer}

% p021_s11
\BookQuestion{一、6}
声音在空气中每秒能传播340米，照这样的速度，实验人员在山谷一侧向对面发射声波，经过6秒收到回声，山谷的宽度是（\hspace{1.8em}）米。\par
\begin{BookAnswer}
6秒内声音走了一个来回，山谷宽度是来回路程的一半：\[\begin{aligned} &340\times6\div2\\={}&2040\div2\\={}&1020\text{(米)}\end{aligned}\]也可先求声音单程传播的时间，再求距离：\[\begin{aligned} &340\times(6\div2)\\={}&340\times3\\={}&1020\text{(米)}\end{aligned}\]检验：\(1020\times2\div340=6\text{(秒)}\)，与听到回声所用时间相符。
\end{BookAnswer}

% p021_s12
\BookQuestion{一、7（每千克价格）}
妈妈用16元买了2千克黄瓜，每千克黄瓜（\hspace{1.8em}）元。\par
\begin{BookAnswer}
\[\begin{aligned} &16\div2\\={}&8\text{(元/千克)}\end{aligned}\]
\end{BookAnswer}

% p021_s13
\BookQuestion{一、7（24元购买质量）}
妈妈用16元买了2千克黄瓜，24元能买（\hspace{1.8em}）千克黄瓜。\par
\begin{BookAnswer}
数量＝总价÷单价：\[\begin{aligned} &24\div8\\={}&3\text{(千克)}\end{aligned}\]
\end{BookAnswer}

\Needspace{8\baselineskip}

\pdfbookmark[2]{原书第22页}{orig-22}

\subsection*{原书第22页}

% p022_s01
\BookQuestion{一、8}
芳芳和悠悠从圆形跑道同一地点，同时出发，相向而行，2分钟后两人相遇，芳芳每分钟跑65米，悠悠每分钟跑85米，这个圆形跑道的周长是（\hspace{1.8em}）米。\par
\begin{BookAnswer}
按题目通常意图，把2分钟理解为首次相遇。两人相向跑2分钟，相遇时共跑一圈：\[\begin{aligned} &(65+85)\times2\\={}&150\times2\\={}&300\text{(米)}\end{aligned}\]若不限定首次，相遇次数未知，便不能仅据此唯一确定周长。
\end{BookAnswer}

% p022_s02
\BookQuestion{一、9}
小华6分钟打432个字，照这样计算，他半小时打（\hspace{1.8em}）个。\par
\begin{BookAnswer}
把打字问题与行程问题对应：每分钟打字数相当于速度，打字总数相当于路程。先统一时间单位：\[\text{半小时}=30\text{分钟}\]先求每分钟打字数，再乘30：\[\begin{aligned} &432\div6\times30\\={}&72\times30\\={}&2160\text{(个)}\end{aligned}\]
\end{BookAnswer}

% p022_s03
\BookQuestion{一、10}
甲、乙两地相距470千米，汽车从甲地出发行驶4小时后，还有30千米才到达终点。这辆汽车每小时行驶（\hspace{1.8em}）千米。\par
\begin{BookAnswer}
实际行驶路程是全程减剩余路程：\[\begin{aligned} &(470-30)\div4\\={}&440\div4\\={}&110\text{(千米/时)}\end{aligned}\]
\end{BookAnswer}

% p022_s04
\BookQuestion{二、1}
卡车以70千米/时的速度从甲地开往乙地，2小时后到达乙地。求甲、乙两地相距多远？列式解答时对应的数量关系式是（\hspace{1.8em}）。\par
\noindent A.速度÷时间=路程\par
\noindent B.总价÷数量=单价\par
\noindent C.速度×时间=路程\par
\noindent D.单价×数量=总价\par
\begin{BookAnswer}
选\(C\)。已知速度和时间，求路程，用“速度×时间＝路程”。本题\(70\times2=140\text{(千米)}\)。
\end{BookAnswer}

% p022_s05
\BookQuestion{二、2}
一支钢笔5元钱，一个文具盒8元钱，买2支钢笔和3个文具盒，共花（\hspace{1.8em}）元。\par
\noindent A.34\par
\noindent B.31\par
\noindent C.30\par
\noindent D.32\par
\begin{BookAnswer}
选\(A\)。2支钢笔和3个文具盒的总价：\[\begin{aligned} &5\times2+8\times3\\={}&10+24\\={}&34\text{(元)}\end{aligned}\]
\end{BookAnswer}

% p022_s06
\BookQuestion{二、3}
北京南站到济南西站的高铁全长约406千米。如右图，G823次列车的平均速度大约是（\hspace{1.8em}）千米/时。\par
\noindent A.230\par
\noindent B.203\par
\noindent C.812\par
\noindent D.820\par
\par\smallskip\Needspace{4\baselineskip}\begingroup\small\setlength{\tabcolsep}{3pt}
\noindent\begin{tabularx}{\linewidth}{|*{4}{>{\centering\arraybackslash}X|}}\hline
出发站 & 出发时刻 & 到达站 & 到达时刻\\\hline
北京南站 & 12:43 & 济南西站 & 14:43\\\hline
\end{tabularx}\par
\noindent 用时2小时\par
\endgroup\smallskip
\begin{BookAnswer}
选\(B\)。从12:43到14:43经过2小时，速度为：\[\begin{aligned} &406\div2\\={}&203\text{(千米/时)}\end{aligned}\]
\end{BookAnswer}

% p022_s07
\BookQuestion{二、4}
下列动物中，奔跑速度最快的是（\hspace{1.8em}）。\par
\noindent A.羚羊每分钟约跑1440米\par
\noindent B.野驴的速度约是1200米/分钟\par
\noindent C.长颈鹿3分钟约跑3千米\par
\noindent D.东北虎1分钟约跑1300米\par
\begin{BookAnswer}
选\(A\)。先把长颈鹿的路程换成米：\[3\text{千米}=3000\text{米}\]统一为米/分：羚羊1440，野驴1200，长颈鹿\(3000\div3=1000\)，东北虎1300。\(1440>1300>1200>1000\)，羚羊最快。
\end{BookAnswer}

% p022_s08
\BookQuestion{二、5}
小华和小明分别从一座桥的两端同时出发。小华的速度是70米/分，小明的速度是60米/分，经过6分钟两人相遇。这座桥长（\hspace{1.8em}）米。\par
\noindent A.390\par
\noindent B.780\par
\noindent C.260\par
\noindent D.130\par
\begin{BookAnswer}
选\(B\)。相向而行，两人路程之和等于桥长：\[\begin{aligned} &(70+60)\times6\\={}&130\times6\\={}&780\text{(米)}\end{aligned}\]
\end{BookAnswer}

% p022_s09
\BookQuestion{二、6}
甲乙两地相距700千米，李叔叔开车从甲地出发，去乙地送货物，每小时行驶78千米，3小时后到达服务区休息，服务区的位置可能在（\hspace{1.8em}）点。\par
\noindent A.①号\par
\noindent B.②号\par
\noindent C.③号\par
\noindent D.④号\par
\BookFigure{figures/p022_s09-01.tex}
\begin{BookAnswer}
选\(B\)。李叔叔3小时行驶\(78\times3=234\)千米，约为700千米全程的三分之一；②号点约在甲至乙全程的三分之一处。图示为位置估计。
\end{BookAnswer}

\Needspace{8\baselineskip}

\pdfbookmark[2]{原书第23页}{orig-23}

\subsection*{原书第23页}

% p023_s01
\BookQuestion{三、1（1）}
直接写得数：\(34\times5=\)。\par
\begin{BookAnswer}
每份34，共5份。\[\begin{aligned} &34\times5\\={}&170\end{aligned}\]
\end{BookAnswer}

% p023_s02
\BookQuestion{三、1（2）}
直接写得数：\(564+39=\)。\par
\begin{BookAnswer}
加39可先加40再减1。\[\begin{aligned} &564+39\\={}&603\end{aligned}\]
\end{BookAnswer}

% p023_s03
\BookQuestion{三、1（3）}
直接写得数：\(105\times21=\)。\par
\begin{BookAnswer}
可拆21为20和1。\[\begin{aligned} &105\times21\\={}&2205\end{aligned}\]
\end{BookAnswer}

% p023_s04
\BookQuestion{三、1（4）}
直接写得数：\(2+8\times75=\)。\par
\begin{BookAnswer}
先乘后加。\[\begin{aligned} &2+8\times75\\={}&602\end{aligned}\]\par\textbf{简便提示（学完运算律后）}\par \[\begin{aligned}&2+8\times75\\={}&2+2\times(4\times75)\\={}&2+2\times300\\={}&2+600\\={}&602\end{aligned}\]
\end{BookAnswer}

% p023_s05
\BookQuestion{三、1（5）}
直接写得数：\(76+0\times2=\)。\par
\begin{BookAnswer}
0乘2得0。\[\begin{aligned} &76+0\times2\\={}&76\end{aligned}\]
\end{BookAnswer}

% p023_s06
\BookQuestion{三、1（6）}
直接写得数：\(703\div19=\)。\par
\begin{BookAnswer}
验算19乘37得703。\[\begin{aligned} &703\div19\\={}&37\end{aligned}\]
\end{BookAnswer}

% p023_s07
\BookQuestion{三、1（7）}
直接写得数：\(24\times18=\)。\par
\begin{BookAnswer}
用乘法分配律，把18拆成20减2，先多算2个24，再减去：\[\begin{aligned} &24\times18\\={}&24\times(20-2)\\={}&24\times20-24\times2\\={}&480-48\\={}&432\end{aligned}\]
\end{BookAnswer}

% p023_s08
\BookQuestion{三、1（8）}
直接写得数：\(327+25=\)。\par
\begin{BookAnswer}
按数位相加。\[\begin{aligned} &327+25\\={}&352\end{aligned}\]
\end{BookAnswer}

% p023_s09
\BookQuestion{三、1（9）}
直接写得数：\(32-8\div4=\)。\par
\begin{BookAnswer}
先除后减。\[\begin{aligned} &32-8\div4\\={}&30\end{aligned}\]
\end{BookAnswer}

% p023_s10
\BookQuestion{三、1（10）}
直接写得数：\(605-37=\)。\par
\begin{BookAnswer}
按减法计算。\[\begin{aligned} &605-37\\={}&568\end{aligned}\]
\end{BookAnswer}

% p023_s11
\BookQuestion{三、2（1）}
计算下列各题：（1）\((540\div27+39)\times17\)。\par
\begin{BookAnswer}
\[\begin{aligned} &(540\div27+39)\times17\\={}&(20+39)\times17\\={}&59\times17\\={}&1003\end{aligned}\]\par\textbf{简便提示（学完运算律后）}\par \[\begin{aligned}&59\times17\\={}&(60-1)\times17\\={}&1020-17\\={}&1003\end{aligned}\]
\end{BookAnswer}

% p023_s12
\BookQuestion{三、2（2）}
计算下列各题：（2）\(75-(416+536)\div34\)。\par
\begin{BookAnswer}
\[\begin{aligned} &75-(416+536)\div34\\={}&75-952\div34\\={}&75-28\\={}&47\end{aligned}\]
\end{BookAnswer}

% p023_s13
\BookQuestion{三、2（3）}
计算下列各题：（3）\(180\times3+728\div8\)。\par
\begin{BookAnswer}
\[\begin{aligned} &180\times3+728\div8\\={}&540+91\\={}&631\end{aligned}\]
\end{BookAnswer}

% p023_s14
\BookQuestion{三、2（4）}
计算下列各题：（4）\(156\div[480\div(276-196)]\)。\par
\begin{BookAnswer}
\[\begin{aligned} &156\div[480\div(276-196)]\\={}&156\div[480\div80]\\={}&156\div6\\={}&26\end{aligned}\]
\end{BookAnswer}

% p023_s15
\BookQuestion{四、1（1）}
（如下图）乐乐和欢欢每天早晨坚持跑步。乐乐每分钟跑220米，欢欢每分钟跑260米。（1）两人同时从家出发，相向而行，3分钟后在小区中心花园相遇，两家相距多少米？\par
\BookFigure{figures/p023_s15-01.tex}
\begin{BookAnswer}
方法一：速度和×时间\par \[\begin{aligned} &(220+260)\times3\\={}&480\times3\\={}&1440\text{(米)}\end{aligned}\]答：两家相距1440米。\par 方法二：分别算两段路程\par \[\begin{aligned} &220\times3+260\times3\\={}&660+780\\={}&1440\text{(米)}\end{aligned}\]答：两家相距1440米。\par
\end{BookAnswer}

% p023_s16
\BookQuestion{四、1（2）}
（接上题）乐乐每分钟跑220米，欢欢每分钟跑260米。（2）两人同时从中心花园出发往学校跑，5分钟后，两人都未到达学校，这时两人相距多少米？（先画图整理条件和问题，再解答。）\par
\BookFigure{figures/p023_s16-01.tex}
\begin{BookAnswer}
两人同地同时同向出发，欢欢较快；路程差就是两人距离。\par 方法一：速度差×时间\par \[\begin{aligned} &(260-220)\times5\\={}&40\times5\\={}&200\text{(米)}\end{aligned}\]答：5分钟后两人相距200米。\par 方法二：分别求路程再相减\par \[\begin{aligned} &260\times5-220\times5\\={}&1300-1100\\={}&200\text{(米)}\end{aligned}\]答：5分钟后两人相距200米。\par
\BookFigure{figures/p023_s16-02.tex}
\end{BookAnswer}

\Needspace{8\baselineskip}

\pdfbookmark[2]{原书第24页}{orig-24}

\subsection*{原书第24页}

% p024_s01
\BookQuestion{四、2}
李叔叔开车从青岛到济南，上午9:00出发，下午1:00到达，汽车平均每小时行驶92千米。青岛到济南的距离是多少千米？\par
\begin{BookAnswer}
先把到达时刻改用二十四时计时法：\[\text{下午}1{:}00=13{:}00\]从9:00到13:00经过\(13-9=4\)小时。\par 主解：速度×总时间\par \[\begin{aligned} &92\times(13-9)\\={}&92\times4\\={}&368\text{(千米)}\end{aligned}\]答：青岛到济南的距离是368千米。\par 分段验算：按中午12时分段\par \[\begin{aligned} &92\times3+92\times1\\={}&276+92\\={}&368\text{(千米)}\end{aligned}\]答：青岛到济南的距离是368千米。\par 两种计时都得到4小时；按上午、下午分段计算用来复核总时间，没有改变“路程＝速度×时间”这一数量关系。
\end{BookAnswer}

% p024_s02
\BookQuestion{四、3}
甲、乙两个工程队分别从两端同时修路，甲队每天修8米，乙队每天修12米，130天后，还有820米没修。这条路一共多少米？\par
\begin{BookAnswer}
先按线段关系整理：全路长＝甲队已修长度＋乙队已修长度＋未修的820米。两队都修130天，既可先求每天合修多少米，也可分别求两队已修多少米。\par 方法一：先合并两队每天修的长度\par \[\begin{aligned} &(8+12)\times130+820\\={}&20\times130+820\\={}&2600+820\\={}&3420\text{(米)}\end{aligned}\]答：这条路一共长3420米。\par 方法二：分别求已修长度再加未修长度\par \[\begin{aligned} &8\times130+12\times130+820\\={}&1040+1560+820\\={}&2600+820\\={}&3420\text{(米)}\end{aligned}\]答：这条路一共长3420米。\par
\end{BookAnswer}

% p024_s03
\BookQuestion{四、4}
某商厦3月份销售248台空调，按照3月份的日销量计算，2026年第2季度能卖多少台空调？\par
\begin{BookAnswer}
3月有31天，4、5、6月分别有30、31、30天，第2季度共91天。\par 方法一：日销量×季度天数\par \[\begin{aligned} &(248\div31)\times(30+31+30)\\={}&8\times(61+30)\\={}&8\times91\\={}&728\text{(台)}\end{aligned}\]答：2026年第2季度能卖728台空调。\par 按3个31天月计算会多算2天，应扣除这2天的销量。\par 方法二：三个月按31天算再调整\par \[\begin{aligned} &248\times3-(248\div31)\times2\\={}&744-8\times2\\={}&744-16\\={}&728\text{(台)}\end{aligned}\]答：2026年第2季度能卖728台空调。\par \par\textbf{简便提示（学完运算律后）}\par 计算季度天数时，两个30天可先合起来凑60，再加31天；这只是合计同样三个月的天数。\[\begin{aligned}&30+31+30\\={}&(30+30)+31\\={}&60+31\\={}&91\text{(天)}\end{aligned}\]
\end{BookAnswer}

% p024_s04
\BookQuestion{四、5}
根据《中华人民共和国道路交通安全法》相关规定，电动自行车在非机动车道行驶时，每分钟不得超过250米。王叔叔点了一份外卖，骑手正在送货，距离目的地2400米，离约定时间还有8分钟。按照规定骑行，骑手能准时送达吗？\par
\begin{BookAnswer}
方法一：比较8分钟最多能走的路程\par \[\begin{aligned} &250\times8\\={}&2000\text{(米)}\end{aligned}\]答：按规定骑行，骑手不能准时送达。\par 因为\(2000<2400\)，8分钟内最多骑行2000米，达不到2400米。\par 方法二：比较准时送达所需的平均速度\par \[\begin{aligned} &2400\div8\\={}&300\text{(米/分)}\end{aligned}\]答：按规定骑行，骑手不能准时送达。\par 因为\(300>250\)，准时送达需要的平均速度超过题目规定。
\end{BookAnswer}
